\subsection{R-25: a ternary bordered-permanent induction and exact size18 dimension}

The first assertion is an exact conditional reduction. A finite
certificate verifies its hypothesis at $m=17$, giving
$\dim\mathcal P_{18,4}=1262$; no uniform vanishing is inferred.

Let $E=\mathbb C^3$, and let $A(x)$ be an $m\times m$ matrix of
linear forms on $E$. Write $P(A)=\operatorname{per}(A)$ and
$G_{ij}(A)=\partial P/\partial A_{ij}$. For vectors $u(x),v(x)$
of linear forms, define
\[
 Q_A(u,v)=DP_A[uv^t]
          =\sum_{i,j}u_iv_jG_{ij}(A),\qquad
 R_A(u,v)=\tfrac12D^2P_A[uv^t,uv^t].
\]
Thus $P(A+tuv^t)=P(A)+tQ_A(u,v)+t^2R_A(u,v)+O(t^3)$.
The derivatives in $u,v$ below mean differentiation in independent
auxiliary coordinates, followed by substitution of the indicated
linear forms. Put
\begin{align*}
 a_i&=\partial_{u_i}Q_A=(Gv)_i,&
 b_j&=\partial_{v_j}Q_A=(G^tu)_j,\\
 t_i&=\partial_{u_i}R_A
          =D^2P_A[e_iv^t,uv^t],&
 t'_j&=\partial_{v_j}R_A
          =D^2P_A[ue_j^t,uv^t].
\end{align*}
All forms in
\begin{equation}\label{border:C}
 \mathcal C_A(u,v)
 =E^*P(A)
   +\sum_i E^*a_i+\sum_j E^*b_j
   +\operatorname{span}\{t_i,t'_j:1\leq i,j\leq m\}
 \ \subseteq\operatorname{Sym}^{m+1}E^*
\end{equation}
have the same degree. This space is generated by $8m+3$ displayed
columns. If at least one of $u,v$ is nonzero, the identity
\begin{equation}\label{border:scaling_relation}
 \sum_i u_i a_i-\sum_jv_jb_j=0
\end{equation}
is a nonzero relation on the corresponding $6m$ linear-multiple
columns. In particular $\dim\mathcal C_A(u,v)\leq8m+2$.

\paragraph{Statement R-25.}
If some ternary data $A,u,v$ satisfy
\begin{equation}\label{border:rank_hypothesis}
 \dim\mathcal C_A(u,v)=8m+2,
\end{equation}
then, with $c_m=\dim\mathcal P_{m,4}$,
\begin{equation}\label{border:increment}
 c_{m+1}\geq c_m+8m+2.
\end{equation}
Consequently, if $m\geq17$ and $c_m=4m^2-2m+2$, the same
equality holds at $m+1$. A proof of
\eqref{border:rank_hypothesis} for all $m\geq17$, together with
R-24, would prove the exact parameter dimension for all those sizes.

\begin{proof}
Adjoin a fourth variable $\ell$, so $V^*=E^*\oplus\mathbb C\ell$.
Extend $A,u,v$ to matrices and vectors of quaternary linear forms.
The condition \eqref{border:rank_hypothesis} on their restrictions
to $\ell=0$ is a nonempty Zariski open condition on these
quaternary parameters. The condition that the differential image
\[
 T_A=\{DP_A[U]:U\in\operatorname{Mat}_m(V^*)\}
\]
have dimension $c_m$ is also nonempty and Zariski open. The
parameter space is irreducible, so both conditions hold
simultaneously for some quaternary $A,u,v$. No maximal-rank
assumption is made on the initial ternary matrix itself.

Consider the single family of quaternary pencils
\[
 B_\varepsilon=
 \begin{pmatrix}A&\varepsilon u\\
                 \varepsilon v^t&\ell\end{pmatrix}.
\]
Every permutation uses either the bottom-right entry, or one
entry from each off-diagonal block. Hence exactly
\begin{equation}\label{border:exact_permanent}
 \operatorname{per}(B_\varepsilon)
       =\ell P(A)+\varepsilon^2Q_A(u,v).
\end{equation}
Differentiation in the bottom-right entry gives every form in
$V^*P(A)$. Differentiation in the upper-right and lower-left
entries, divided by $\varepsilon$, gives every form in
$\sum V^*a_i+\sum V^*b_j$, for $\varepsilon\ne0$.

The remaining $2m$ columns come from exact cancellations, not
from independently chosen deformations. Fix a constant vector
$h\in\mathbb C^m$. Vary the upper-right block by $\ell h$ and
the upper-left block by $-\varepsilon hv^t$, leaving all other
entries fixed. The differential of
\eqref{border:exact_permanent} in that direction is
\begin{align*}
 \varepsilon\ell\,DP_A[hv^t]
  -\varepsilon\ell\,DP_A[hv^t]
  -\varepsilon^3D_AQ_A(u,v)[hv^t]
   &=-\varepsilon^3D^2P_A[hv^t,uv^t]\\
   &=-\varepsilon^3\sum_i h_it_i.
\end{align*}
Every entry of this parameter variation is a linear form, as
required. Varying instead the lower-left block by $\ell h^t$
and the upper-left block by $-\varepsilon uh^t$ similarly gives
$-\varepsilon^3\sum_jh_jt'_j$. Division by
$-\varepsilon^3$ is permitted in the differential image for
every nonzero $\varepsilon$.

Choose $8m+2$ of the fixed polynomial columns in
\eqref{border:C} whose restrictions to $\ell=0$ are independent.
All belong to the differential image of this same
$B_\varepsilon$, for every $\varepsilon\ne0$.
Choose constant quaternary variations $U_1,\ldots,U_{c_m}$ such
that $DP_A[U_i]$ form a basis of $T_A$. Their derivatives at
$B_\varepsilon$ are
\[
 \ell DP_A[U_i]+\varepsilon^2D_AQ_A(u,v)[U_i].
\]
At $\varepsilon=0$ these become a basis of $\ell T_A$.
The combined $c_m+8m+2$ limit columns are independent: restriction
to $\ell=0$ first annihilates all coefficients of the selected
normal columns, and multiplication by the nonzero polynomial
$\ell$ is injective on $T_A$.

A coefficient minor of the combined columns is therefore a
polynomial in $\varepsilon$ with nonzero constant term. It is
nonzero at some $\varepsilon\ne0$, proving
\eqref{border:increment}. Finally the torus upper bound increases
by exactly
\[
 [4(m+1)^2-2(m+1)+2]-[4m^2-2m+2]=8m+2,
\]
so equality propagates whenever the hypothesis is available.
\end{proof}

\paragraph{Exact size18 corollary.}
The rank hypothesis \eqref{border:rank_hypothesis} holds at $m=17$.
Consequently
\begin{equation}\label{R25:size18}
 c_{18}=1262=4\cdot18^2-2\cdot18+2,
 \qquad \epsilon_{18}=0.
\end{equation}
The genuinely padded pairs $(18,n)$, $19\leq n\leq25$, therefore
satisfy the strict four-variable dimension comparison. In
particular $(18,25)$ admits the G-17/G-18 stretched direction
$\lambda=(46r,18r,18r,18r)$ in coordinate degree $4r$; G-12
excludes its rectangular sector.

\begin{proof}
Here is the complete finite specification of the ternary
certificate. All indices are zero-based. Set $p=1021$ and
$s_0=20330804$, and generate
\[
 s_{a+1}=1664525s_a+1013904223\pmod{2^{32}},
 \qquad c_a=\lfloor s_{a+1}/2^8\rfloor\bmod1021.
\]
Use the first $3\cdot17^2=867$ coefficients, in row, column,
variable order, for $A$; use the next $51$ for $u$ and the final
$51$ for $v$, in vector-index, variable order. These are integer
linear forms in $(x_0,x_1,x_2)$.

List the $190=\binom{20}{2}$ nonnegative integer triples of sum
$18$ in decreasing lexicographic order. Evaluate the following
$139$ columns at those points, in precisely the indicated order:
\begin{equation}\label{R25:columns}
 \begin{array}{ll}
  x_kP(A) &(0\leq k<3),\\
  x_ka_i &(0\leq i<17,\ 0\leq k<3),\\
  x_kb_j &(0\leq j<17,\ 0\leq k<3),\\
  t_i &(0\leq i<17),\\
  t'_j &(0\leq j<17).
 \end{array}
\end{equation}
The $138\times138$ minor on the first $138$ rows and all columns
except column $104$ is, in increasing row and column order,
\begin{equation}\label{R25:minor}
 \det C_*=501\pmod{1021}.
\end{equation}
For completeness the polynomial recurrence determining these
entries without a symbolic permanent expansion is as follows.
At each numerical point put $H=uv^t$. For a subset $S$ of
columns and $i=|S|-1$, recursively compute the polynomial
$F_S(t)$ modulo $t^3$ by
\[
 F_\varnothing(t)=1,\qquad
 F_S(t)=\sum_{j\in S}(A_{ij}+tH_{ij})F_{S\setminus\{j\}}(t)
                         \pmod{t^3}.
\]
The constant, linear and quadratic coefficients of $F_{[17]}$
are exactly $P,Q,R$. Differentiating this finite recurrence in
the auxiliary coordinates $u,v$ gives all columns in
\eqref{R25:columns}. The source
\texttt{check\_border\_ternary.cpp} implements this recurrence
and reverse differentiation, then exact elimination on the
specified evaluation matrix, yielding \eqref{R25:minor}.

A distinct implementation, \texttt{check\_border\_ryser.cpp},
recomputed all $26{,}410$ entries from Gray-code Ryser sums:
prefix and suffix products of the row-sum linear factors give
the degree-zero and degree-one complementary products, and
explicit auxiliary differentiation gives the derivatives of
$Q,R$. Every entry agreed. An independent agent reconstructed
all generator coefficients and points, verified the complete
scaling relation \eqref{border:scaling_relation}, and checked
\eqref{R25:minor} by modular fraction-free Bareiss elimination.
The latter check concerns the minor and scaling relation; it
is separate from the author's all-entry Ryser recomputation.

The evaluated columns are evaluations of the homogeneous forms
in \eqref{border:C}. Nonzero determinant modulo $1021$ gives
rank at least $138$ over $\mathbb C$. The universal relation
\eqref{border:scaling_relation} gives the opposite bound, so
\eqref{border:rank_hypothesis} holds at $m=17$.
R-24 and \eqref{border:increment} now give
$c_{18}\geq1124+138=1262$. The torus bound gives equality, and
R-23 identifies the vanishing homological defect.

Finally $c_{18}+3=1265>2\cdot25^2+2=1252$, proving the padded
range by R-14. At $(m,n)=(18,25)$, $a=7$ and
$\nu=(46,18,18,18)$ has $|\nu|=100=4n$, $\nu_4>0$, and
$\nu_1>4a$. Thus G-17/G-18 apply. Since $25>4\cdot18/3$,
G-12 excludes rectangles.
\end{proof}

The uniform gate left by this result is a proof of
\eqref{border:rank_hypothesis} for every $m\geq18$, or a
replacement induction giving the same dimension increment.
The finite certificate alone establishes neither assertion.
R-27 below closes this gate by a uniform endpoint-jet argument,
proving the hypothesis for every $m\geq17$ and hence exact
dimension at every subsequent size.
